\chapter{Az RWLang CLI és a napi fejlesztői/üzemeltetői workflow}

Az RWLang két telepített programot ad: a hosszú életű \code{rwlang-server} processzt és a rövid életű \code{rwlang-cli} operátori eszközt. A kettőt érdemes következetesen külön kezelni. A szerver HTTP-t szolgál ki és runtime policyt érvényesít; a CLI forrást ellenőriz, migrációt kezel és local-auth adminisztrációt végez.

Ez a fejezet nem új nyelvi feature-öket vezet be. Azt mutatja meg, hogyan áll össze a napi munka egy ismételhető parancsfolyamattá a szerkesztéstől a production ellenőrzésig.

\section{Melyik feladathoz melyik program?}

\begin{longtable}{L{0.29\textwidth}L{0.23\textwidth}L{0.39\textwidth}}
\toprule
Feladat & Eszköz & Megjegyzés \\
\midrule
RWLang forrás ellenőrzése & \code{rwlang-cli check} & Compiler-szintű, nem indít HTTP szervert. \\
Migration status/verify/apply & \code{rwlang-cli migrate} & DB URL secret fájlból. \\
Local-auth store admin & \code{rwlang-cli auth} & SQLite identity store, operátori művelet. \\
Server config preflight & \code{rwlang-server} & \code{--check-config}: config + app + lokális policy/resource fájlok ellenőrzése. \\
Effective config megtekintése & \code{rwlang-server} & \code{--print-effective-config}; a secret értékek redaktálva maradnak. \\
Webalkalmazás futtatása & \code{rwlang-server} & Hosszú életű processz, proxy vagy közvetlen HTTPS mögött. \\
Workspace release build & \code{cargo build} & \code{--locked --release} módban készül a két bináris. \\
Platform teljes verifikációja & \code{./verify.sh} & Az RWLang repository release-evidence része, nem minden alkalmazás saját parancsa. \\
\bottomrule
\end{longtable}

A production példákban a programok helye:

\begin{lstlisting}
/usr/local/bin/rwlang-cli
/usr/local/bin/rwlang-server
\end{lstlisting}

A shell normál \code{PATH} beállítása mellett ezért a további példákban elég a rövid programnevet használni.

\section{A CLI nem a szerver alternatívája}

A \code{rwlang-cli} nem indít listenert, nem szolgál ki route-ot és nem tart fenn session state-et. Egy parancsot végrehajt, majd kilép. Ez biztonsági szempontból is hasznos: migration vagy user-admin művelethez nem kell külön admin HTTP endpointot nyitni.

A másik irányban a \code{rwlang-server} nem futtat automatikus migrationt startupkor. A schema változás explicit operátori döntés marad.

\section{Az első parancs: \texttt{rwlang-cli check}}

Napi fejlesztésben minden érdemi forrásmódosítás után ez legyen az első kapu:

\begin{lstlisting}
rwlang-cli check main.rw
\end{lstlisting}

Siker esetén a CLI rövid összegzést ír a fordító által látott model-, query-, page-, action- és route-számról. Hiba esetén nem érdemes szervert indítani: előbb a compile-time contractot kell kijavítani.

A \code{check} többek között olyan hibákat fog meg, mint:

\begin{itemize}
  \item ismeretlen típus vagy hibás nyelvi szintaxis;
  \item query projection és model shape eltérése;
  \item hibás typed bind/decode contract;
  \item route/policy vagy resource profile hivatkozási hiba, amely compiler-szinten bizonyítható;
  \item tiltott raw HTML/SQL jellegű konstrukció.
\end{itemize}

A \code{check} viszont nem bizonyítja, hogy a DB vagy Redis hálózaton elérhető, és nem helyettesíti a HTTP integration tesztet.

\section{Forrásfából futtatás és telepített CLI}

Az RWLang repository fejlesztésekor ugyanaz a program Cargo-n keresztül is futtatható:

\begin{lstlisting}
cargo run --locked -p rwlang-cli -- check main.rw
\end{lstlisting}

Ez fejlesztői kényelmi forma. Production runbookban az immutable, telepített release-binárist használd:

\begin{lstlisting}
/usr/local/bin/rwlang-cli check /srv/rwlang/app/main.rw
\end{lstlisting}

Ne keverd a kettőt egy deploymentben. A production művelet legyen egyértelműen egy release artifacthoz köthető.

\section{Migration workflow}

A migration parancs három műveletet ad:

\begin{description}
  \item[\code{status}] megmutatja az applied és pending migrationöket;
  \item[\code{verify}] ellenőrzi a migration historyt és checksum-contractot;
  \item[\code{apply}] lock mellett alkalmazza a pending migrationöket.
\end{description}

A DB URL nem parancssori plaintext secret, hanem egy egysoros secret fájl:

\begin{lstlisting}
/run/secrets/rwlang/db-url
\end{lstlisting}

A tipikus ellenőrzési sorrend:

\begin{lstlisting}
rwlang-cli migrate status \
  --dir migrations \
  --db-url-file /run/secrets/rwlang/db-url

rwlang-cli migrate verify \
  --dir migrations \
  --db-url-file /run/secrets/rwlang/db-url
\end{lstlisting}

Jóváhagyott schema-változásnál:

\begin{lstlisting}
rwlang-cli migrate apply \
  --dir migrations \
  --db-url-file /run/secrets/rwlang/db-url \
  --lock-timeout-secs 30

rwlang-cli migrate verify \
  --dir migrations \
  --db-url-file /run/secrets/rwlang/db-url
\end{lstlisting}

A \code{--lock-timeout-secs} értéke pozitív egész. Ha a migration lock nem szerezhető meg a megadott időn belül, a művelet hibával álljon meg; ne legyen párhuzamos schema-módosítás.

\subsection*{Az \texttt{--allow-insecure-db} nem kényelmi kapcsoló productionre}

A migration CLI is ugyanazt a secure-by-default szemléletet követi, mint a szerver. Az \code{--allow-insecure-db} csak tudatos, kontrollált környezetben használható. Productionben a megfelelő TLS/DB security legyen az alapértelmezett, ne egy állandósult bypass flag.

\section{Local-auth adminisztráció}

Ha az alkalmazás local authot használ, az identity store kezelése szintén a CLI feladata. Először inicializálni kell a store-t:

\begin{lstlisting}
rwlang-cli auth init \
  --db-url-file /run/secrets/rwlang/local-auth-db-url
\end{lstlisting}

Felhasználó létrehozása:

\begin{lstlisting}
rwlang-cli auth user-add \
  --db-url-file /run/secrets/rwlang/local-auth-db-url \
  --username alice \
  --password-file /run/secrets/rwlang/bootstrap-password \
  --role Editor
\end{lstlisting}

Ha nincs \code{--role}, új usernél a CLI a \code{User} role-t használja. Több role több \code{--role} kapcsolóval adható meg.

Jelszócsere és account állapot:

\begin{lstlisting}
rwlang-cli auth password-set \
  --db-url-file /run/secrets/rwlang/local-auth-db-url \
  --username alice \
  --password-file /run/secrets/rwlang/new-password

rwlang-cli auth disable \
  --db-url-file /run/secrets/rwlang/local-auth-db-url \
  --username alice

rwlang-cli auth enable \
  --db-url-file /run/secrets/rwlang/local-auth-db-url \
  --username alice
\end{lstlisting}

Role-csere:

\begin{lstlisting}
rwlang-cli auth roles-set \
  --db-url-file /run/secrets/rwlang/local-auth-db-url \
  --username alice \
  --role Editor --role Publisher
\end{lstlisting}

A \code{roles-set} a megadott role-listát tekinti az új állapotnak; az operátor ne úgy kezelje, mint egyetlen role inkrementális hozzáadását.

\section{TOTP adminisztráció különösen érzékeny művelet}

TOTP enrolment:

\begin{lstlisting}
rwlang-cli auth totp-enroll \
  --db-url-file /run/secrets/rwlang/local-auth-db-url \
  --username alice \
  --recovery-count 8
\end{lstlisting}

A parancs kiírja a TOTP secretet, az \code{otpauth} URI-t és az egyszer megjelenő recovery code-okat. Emiatt ezt ne futtasd CI logban, shell transcriptben vagy közös terminálon. Az operátori folyamatnak biztosítania kell a secret és recovery code-ok azonnali, biztonságos átadását és tárolását.

TOTP kikapcsolása:

\begin{lstlisting}
rwlang-cli auth totp-disable \
  --db-url-file /run/secrets/rwlang/local-auth-db-url \
  --username alice
\end{lstlisting}

\section{A secret fájl is input contract}

A CLI a DB URL és password fájlokat nem korlátlanul olvassa. A jelenlegi contract szerint a secret URL fájl:

\begin{itemize}
  \item regular file;
  \item nem symlink;
  \item legfeljebb 16 KiB;
  \item pontosan egy nem üres sort tartalmaz.
\end{itemize}

A password file szintén regular, nem symlink fájl; legfeljebb 4096 byte, egyetlen sor, a jelszó pedig 12--1024 byte közötti.

Ezért ne process substitutionnel, symlinkkel vagy több soros secret template-tel próbáld megkerülni a contractot. A CLI boundary szándékosan szűk.

\section{Server preflight: a CLI után a szerver saját ellenőrzése jön}

A forrás és migration ellenőrzése után a production konfigurációt a szerverrel validáld:

\begin{lstlisting}
rwlang-server \
  --config /usr/local/etc/rwlang/server.toml \
  --check-config
\end{lstlisting}

Ez beolvassa a trusted TOML-t, lefordítja az alkalmazást, ellenőrzi a route-rate/resource policy hivatkozásokat, TLS konfigurációt, valamint a deklarált storage/static rootok lokális feltételeit.

Ha a precedence-t vagy az effektív policyt kell auditálni:

\begin{lstlisting}
rwlang-server \
  --config /usr/local/etc/rwlang/server.toml \
  --print-effective-config
\end{lstlisting}

A kimenet kifejezetten secret-redacted. Ettől még operátori információ, ezért automation logban csak indokolt retentionnel tárold.

\section{Napi fejlesztői ciklus}

Egy kis vagy közepes alkalmazásnál jó alapfolyamat:

\begin{lstlisting}
edit .rw source
-> rwlang-cli check main.rw
-> rwlang-cli migrate verify
-> if schema changed: migrate status/apply/verify on dev DB
-> rwlang-server --config server.dev.toml --check-config
-> start rwlang-server
-> HTTP/form/JSON smoke tests
-> automated application tests
\end{lstlisting}

A sorrend lényege, hogy az olcsóbb és determinisztikusabb hibákat előbb fogjuk meg. Nincs értelme HTTP smoke tesztet futtatni olyan forrásra, amely már compiler-szinten hibás.

\section{Application release workflow}

Egy production release minimális operátori folyamata:

\begin{lstlisting}
immutable release selected
-> rwlang-cli check deployed main.rw
-> rwlang-server --check-config
-> rwlang-cli migrate status
-> rwlang-cli migrate verify
-> backup/recovery gate
-> approved rwlang-cli migrate apply, if pending
-> rwlang-cli migrate verify
-> controlled service restart/switch
-> /health/live
-> /health/ready
-> application smoke test
-> logs + metrics + audit review
\end{lstlisting}

A migration \code{apply} az egyetlen mutáló schema-lépés ebben a folyamatban; külön approval pontként kezeld.

\section{Az RWLang platform release workflowja más}

Ha magát az RWLang repositoryt buildeljük, a Rust workspace release-evidence is része a folyamatnak:

\begin{lstlisting}
./verify.sh
cargo build --locked --release \
  -p rwlang-server -p rwlang-cli
\end{lstlisting}

A kész binárisok:

\begin{lstlisting}
target/release/rwlang-server
target/release/rwlang-cli
\end{lstlisting}

Telepítés:

\begin{lstlisting}
sudo install -m 0755 target/release/rwlang-server \
  /usr/local/bin/rwlang-server
sudo install -m 0755 target/release/rwlang-cli \
  /usr/local/bin/rwlang-cli
\end{lstlisting}

A saját RWLang webalkalmazásod napi workflowjához nem kell minden alkalommal a nyelv teljes Rust workspace-ét verifikálni. Ott a saját source-, migration-, config- és integration tesztjeid a relevánsak.

\section{Automatizálás: read-only és mutáló lépések legyenek külön}

CI/CD-ben különítsd el a biztonságosan ismételhető ellenőrzéseket a state-et módosító műveletektől.

\textbf{Read-only/preflight:}

\begin{lstlisting}
rwlang-cli check main.rw
rwlang-cli migrate status \
  --dir migrations --db-url-file /run/secrets/rwlang/db-url
rwlang-cli migrate verify \
  --dir migrations --db-url-file /run/secrets/rwlang/db-url
rwlang-server --config /usr/local/etc/rwlang/server.toml \
  --check-config
rwlang-server --config /usr/local/etc/rwlang/server.toml \
  --print-effective-config
health/readiness GET
\end{lstlisting}

\textbf{Mutáló/operator-approved:}

\begin{lstlisting}
rwlang-cli migrate apply \
  --dir migrations --db-url-file /run/secrets/rwlang/db-url
rwlang-cli auth init \
  --db-url-file /run/secrets/rwlang/local-auth-db-url
# A tobbi auth muvelet ugyanilyen explicit --db-url-file
# mellett kapja a sajat --username/--password-file/--role opcioit.
\end{lstlisting}

Az auth-admin parancsokat általában ne tedd automatikus application deployment pipeline-ba. Ezek identity-admin műveletek, külön audit- és hozzáférési boundaryval.

\section{Gyakori hibák}

\begin{description}
  \item[\code{check} sikeres, a server mégsem indul] A compile-time contract rendben van, de konfigurációs, TLS-, filesystem- vagy policy-startup hiba van. Futtasd a \code{--check-config} ellenőrzést és olvasd a server logot.
  \item[\code{--check-config} sikeres, readiness piros] A lokális config helyes, de runtime dependency -- például DB vagy Redis -- nem érhető el vagy nem egészséges.
  \item[\code{migrate verify} hibázik] Ne futtass reflexből \code{apply}-t. Előbb tisztázd a history/checksum eltérést.
  \item[Auth CLI nem nyitja meg a DB-t] Ellenőrizd, hogy local-auth SQLite URL-t adtál-e, és a store inicializálva van-e.
  \item[Secret file hibát kap] Ne lazítsd a fájlcontractot; regular, nem symlink, egysoros secretet adj.
\end{description}

\section{Egy oldalnyi operátori runbook}

Ha csak a minimális sorrendet akarod megjegyezni:

\begin{lstlisting}
# source
rwlang-cli check main.rw

# schema
rwlang-cli migrate status ...
rwlang-cli migrate verify ...

# config
rwlang-server --config /usr/local/etc/rwlang/server.toml \
  --check-config

# approved change
rwlang-cli migrate apply ...
rwlang-cli migrate verify ...

# service
systemctl restart rwlang

# end-to-end health through the public TLS edge
curl -fsS https://www.example.com/health/live
curl -fsS https://www.example.com/health/ready

# proxy topology: optional local upstream probe
curl -fsS http://127.0.0.1:8080/health/ready

# evidence
review server/access/audit log + metrics
\end{lstlisting}

A konkrét systemd unit neve deploymentenként eltérhet; a könyvben szereplő \code{rwlang} csak példa. A fontos invariáns a sorrend: \textbf{source → schema → config → approved mutation → service → health → evidence}.
